\renewcommand{\thechapterdisplay}{Vigésimo sétimo dia}
\renewcommand{\thechapterrunner}{Oitavo modo de aliviar as almas do Purgatório: o ato heroico de caridade}
\chapter[Vigésimo sétimo dia --- Oitavo modo de aliviar as almas do Purgatório: o ato heroico de caridade]{Oitavo modo de aliviar as almas do Purgatório: o ato heroico de caridade}
\section*{Sua natureza}

O ato heroico de caridade em favor das almas do Purgatório é uma doação voluntária do aspecto satisfatório de nossas obras, durante nossa vida, e dos sufrágios que nos serão aplicados após a morte. Nós os depositamos nas mãos de Maria, afim de que essa terna mãe os distribua como for de seu agrado, às almas que ela deseja libertar de suas penas. Em virtude desse ato generoso, feito uma única vez e válido para sempre, para o tempo e para a eternidade, nós nos despojamos somente de nossas obras satisfatórias, que têm por objetivo quitar perante Deus as penas que lhe devemos por nossos pecados, e dos sufrágios que a piedade dos vivos nos oferecerá após nossa morte. Mas conservamos o mérito de nossas ações e seus frutos impetratórios, que alcança de Deus as graças. Por isso, essa doação não impede os padres de oferecerem o Santo Sacrifício por aqueles que lhes pedem, nem os fiéis de rezarem por eles e por seus parentes. Este voto, se for feito, não obriga sob pena de pecado, e pode-se mudar de intenção assim que se desejar. Chama-se \textit{voto heroico}, porque aquele que o faz de bom coração realiza um ato prodigioso de abnegação: ele parece esquecer-se de si mesmo e sacrificar-se para socorrer aqueles que ele ama até o heroísmo.

Não há fórmula prescrita. É suficiente fazer essa oferta de coração, ou de exprimí-la nestes termos: \textit{Ó Maria, Consoladora dos aflitos, eu vos ofereço e vos dou todas as satisfações associadas às minha boas obras, todos os sufrágios que me acompanharão para além do túmulo. Distribuí-os a vossos filhos desafortunados do Purgatório, particularmente às almas de meus pais, dos meus amigos, de meus benfeitores. Assim seja.}


\section*{Suas vantagens}

Primeiramente, essa prática é muito útil às almas do Purgatório. Que socorro não recebem todos os dias, a cada instante do dia, de todas nossas obras satisfatórias, e sobretudo daquelas que nos serão aplicadas durante nossa vida, à nossa morte, e após nossa passagem para a eternidade. É um doce e contínuo orvalho de sufrágios e de indulgências que caem sem interrupção sobre as almas abrasadas no lugar de expiação e apaga o fogo que as devora. 

Em segundo lugar, essa doação heroica não é menos vantajosa para nós. Deus, que é tão bom, tão liberal, não nos dará o cêntuplo de tudo o que fizermos por esses filhos infelizes? \textit{Dai}, diz Ele (Lucas 6:38), \textit{e vos será dado, e recebereis uma boa medida, calcada, sacudida, transbordante; porque com a mesma medida com que medirdes os outros, vós também sereis medidos.} A Santíssima Virgem, a quem teremos confiado todos os nossos tesouros espirituais para aliviar seus irmãos e os nossos, não virá ela em nosso socorro? Esse abandono filial não nos dá direito às liberalidades de sua misericórdia? Emfim, não poderemos contar com a gratidão das almas que tivermos aliviado? Assim, é uma crença geral que aquele que fez o voto heroico não tem muito por que temer o Purgatório. Deus lhe fornecerá o modo de evitá-lo, ou ao menos, não deixará que sofra muito tempo em suas chamas. 

Eu vos aconselho, alma cristã, de fazer hoje mesmo o voto de caridade heroico, se não o tiverdes feito ainda. Segui o exemplo de um número muito considerável de personagens ilustres em dignidade, em ciência, em santidade. Segui os conselhos do venerado Pio IX, que recomendava frequentemente o voto heroico e que o enriqueceu com indulgências. 


\section*{Exemplo}

Lemos na vida de Santa Gertrudes que, desde sua mais tenra idade, ela aprendeu a oferecer todas as orações e todas as suas boas obras pela intenção das almas do Purgatório, pelo voto heroico de caridade. Essa prática era tão agradável a Deus, que frequentemente o divino Salvador comprazia-se em lhe encaminhar as Almas mais necessitadas, e estas, libertadas por sua piedosa caridade, mostravam-se depois a ela em meio à glória, e a agradeciam com efusão, prometendo-lhe não esquecê-la no Céu. Gertrudes tinha passado sua vida nesse santo exercício e, cheia de confiança, ela via com calma a morte se aproximar, quando o inimigo infernal veio dizer-lhe que ela havia se despojado de todo o mérito satisfatório de suas boas obras, e que ela iria cair no Purgatório, para ali expiar todas as suas faltas, em longos sofrimentos. Esse tormento do espírito a lançou numa tal desolação, que seu celeste esposo, Nosso Senhor, se dignou vir para a consolar. \textit{Por que, ó Gertrudes, estás triste, tu que há pouco gozavas de uma perfeita serenidade?} --- \textit{Ah! Senhor, em que deplorável situação eu me encontro! Eis que a morte se aproxima, e estou privada da satisfação de minhas boas obras que apliquei às almas do Purgatório: como poderei agora pagar as dívidas que contraí eu mesma para com a vossa justiça?} --- O salvador respondeu então com carinho. \textit{Não temas, minha bem-amada, pois tu tens, ao contrário, por tua caridade para com os mortos, aumentado a soma de teus méritos satisfatórios, e não somente possuis o bastante para expiar tuas leves faltas, mas adquiriste um grau muito alto de glória na beatitude eterna. É assim que minha clemência reconhecerá, por uma generosa recompensa, tua devoção pelos mortos; e tu virás em breve ao Paraíso receber o cêntuplo de tudo o que fizeste por eles.} Como são encorajadoras essas palavras do divino Mestre!


\textbf{Oremos.} Eu vos suplico, ó meu Deus, de aprovar e de confirmar este voto que ofereço pela Vossa glória e pela salvação de minha alma. Eu o ofereço ainda para pagar todas as dívidas das almas do Purgatório que a Santa Virgem quer libertar. Tomo como testemunhas de meu compromisso todos os eleitos da Igreja militante na terra e padecente no Purgatório. Nesta consideração, sede propício, Senhor, àqueles que sofrem no lugar de expiação. Que descansem em paz!


Em seguida, rezar as orações do final deste livro.
